% Template for computer science thesis at the TU Munich
% Authors Benedikt Mas y Parareda, Johannes Becker, Gunnar Schroeder, Elmar Juergens

\documentclass[11pt,a4paper]{article}

%Enable DVI forward search
%\usepackage[active]{srcltx}

\usepackage{multicol}
\usepackage[utf8x]{inputenc}
\usepackage[T1]{fontenc}
\usepackage{amsmath,amssymb,amsfonts}
\usepackage{color}
\usepackage{graphicx}
\usepackage{rotating}
\usepackage{listings}
\usepackage{cite}
\usepackage{url}
\usepackage{latexsym}
\usepackage{makeidx}
\usepackage{color}
\usepackage{natbib}
\usepackage{array}
\usepackage{todonotes}

\lstset{
        basicstyle=\scriptsize\ttfamily, % Standardschrift
        numbers=left,               % Ort der Zeilennummern
        numberstyle=\tiny,          % Stil der Zeilennummern
        %stepnumber=2,               % Abstand zwischen den Zeilennummern
        numbersep=5pt,              % Abstand der Nummern zum Text
        tabsize=2,                  % Groesse von Tabs
        %extendedchars=true,         %
        breaklines=true,            % Zeilen werden Umgebrochen
        keywordstyle=\color{red}\ttfamily,
   		frame=b,         
%        keywordstyle=[1]\textbf,    % Stil der Keywords
 %       keywordstyle=[2]\textbf,    %
  %      keywordstyle=[3]\textbf,    %
   %     keywordstyle=[4]\textbf,   \sqrt{\sqrt{}} %
        stringstyle=\color{blue}\ttfamily, % Farbe der String
        showspaces=false,           % Leerzeichen anzeigen ?
        showtabs=false,             % Tabs anzeigen ?
        xleftmargin=17pt,
        framexleftmargin=17pt,
        framexrightmargin=5pt,
        framexbottommargin=4pt,
        %backgroundcolor=\color{lightgray},
        %showstringspaces=false,      % Leerzeichen in Strings anzeigen ?        
		language=Java
}
\lstloadlanguages{Java}
\usepackage{caption}
\DeclareCaptionFont{white}{\color{white}}
\DeclareCaptionFormat{listing}{\colorbox[cmyk]{0.88, 0.44, 0.0,0.0}{\parbox{\textwidth}{\hspace{15pt}#1#2#3}}  }
\captionsetup[lstlisting]{format=listing,labelfont=white,textfont=white, singlelinecheck=false, margin=0pt, font={bf,footnotesize}}

\definecolor{darkgreen}{cmyk}{0.7, 0, 1, 0.5}
\definecolor{darkblue}{rgb}{0.1, 0.1, 0.5}

\lstdefinelanguage{diff}
{
keywords={+, -, \ , @@, diff, index, new},
sensitive=false,
morecomment=[l][""]{\ },
morecomment=[l][\color{darkgreen}]{+},
morecomment=[l][\color{red}]{-},
morecomment=[l][\color{darkblue}]{@@},
morecomment=[l][\color{darkblue}]{diff},
morecomment=[l][\color{darkblue}]{index},
morecomment=[l][\color{darkblue}]{new},
morecomment=[l][\color{darkblue}]{similarity},
morecomment=[l][\color{darkblue}]{rename},
}

\makeindex


\usepackage{geometry,mflogo,xspace,texnames,path,booktabs,bm}
\usepackage[hyperindex,bookmarks,pdfborder=0,plainpages=false,pdfpagelabels]{hyperref}


%Settings applicable to the complete document

%Breaks URLs properly and draws them in a nice font
\urlstyle{sm}

\renewcommand{\ttdefault}{pcr} % Courier has a bold shape, default tt does not


%Decrease the default indentation of paragraphs
\parindent=0.3cm

%new or changed commands
\renewcommand\contentsname{Table of Content}
\newcommand{\chapref}[1]{Chapter~\ref{#1}}
\newcommand{\secref}[1]{Section~\ref{#1}}
\newcommand{\appref}[1]{Appendix~\ref{#1}}
\newcommand{\tabref}[1]{Table~\ref{#1}}
\newcommand{\figref}[2][]{Figure~\ref{#2}#1}
\newcommand{\listref}[2][]{Listing~\ref{#2}#1}

\newcommand{\UH}{\textit{Unknown Horizons}}
\newcommand{\OS}{open-source}
\newcommand{\BOW}{\textit{Battle for Wesnoth}}
\newcommand{\AD}{\textit{0 A.D.}}
\newcommand{\GLEST}{\textit{Mega Glest}}

\begin{document}
%\bibliographystyle{plainnat}

%Title Page
\pagestyle{empty}

\begin{titlepage}
\begin{center}

\begin{figure}[!htb]
	\begin{center}
		\includegraphics[scale=0.50]{pics/ntnu}
	\end{center}
\end{figure}
\begin{figure}[!htb]
	\begin{center}
		\includegraphics[scale=0.40]{pics/idi}
	\end{center}
\end{figure}
\begin{LARGE}
\vspace{1.2in}

\end{LARGE}
\begin{Huge}
Report: Poker AI Implementation in Java
\vspace{1.7in}


\end{Huge}
\begin{LARGE}
Henning Funke \& Thomas Kinnen \vspace{0.6in}
\end{LARGE}

\begin{Large}
IT3105 - Kunstig intelligens programmering
\vspace{0.8in}
\end{Large}



\end{center}

\end{titlepage}

%\input{abstract.tex}
%\input{acknowledgement.tex}


% Table of Contents
\setcounter{tocdepth}{1}                % Sets depth of table of contents. 0 is chapter, 1 is sections, 2 is subsections
\setcounter{secnumdepth}{2}             % Sets depth of numbering of toc contents
\tableofcontents

\pagestyle{headings}

% Chapters. Each one in its own file

\section{Implementations for parts 1, 2 and 3}

In the first three parts we implement several systems that determine if 
a hypothesis is entailed by a text.


\subsection{Basic framework}

We created a framework that enables us to create, use and evaluate
a variety of RTE-systems in an easy way. So we have all systems of all parts
of the project embedded in one big test environment.
It loads the syntactically parsed data and wraps it in an object oriented structure 
\textit{Text}. It is able to provide the data either as a linked graph or as a 
list of words.

To evaluate the performance of our systems we decided to implement our own evaluation
code. The huge amount of file operations, that goes along with using the provided 
python script has a big impact on the execution time.
To achieve a common basis for evaluation and easy implementation of new systems we make use of 
object oriented techniques. All of our RTE-systems are derived from the interface
\textit{IEntailmentRecognizer}, that requires functionality to estimate the entailment
propability of two given \textit{Text} instances.

Further we implemented an abstract class \textit{BasicMatcher}, which is the basis for a
category of RTE-systems, that determine the entailment by comparing the sentences word by word.
Thus we are able to create new systems of that category by implementing only one procedure, 
that chooses or calculates the desired property.

\subsection{Systems in part 1}
\subsection{Systems in part 2}
\subsection{Systems in part 3}



\section{Building our own RTE system}
In this section we describe our best RTE system.

\subsection{Considerations from the prior systems}
All prior systems turned out to be worse than the normal lemma matching, Which to this point was still our  best system
with a score of 63.25\%. We did not see any big potential in improving the basic lemma matching further, so we decided
to try to tune our machine learning matcher \textit{BasicMahoutMatcher} in the form of \textit{MahoutMatcher}. Our basic
idea was to use the best non-machine-learner systems as features for the machine learner, hoping this would give the
machine learner a good basis to work with.

\subsection{Implementation}
We already had the basic \textit{BasicMahoutMatcher} from part three, which we decided to tune. The base of this matcher
was the \textit{OnlineLogisticRegression} learning algorithm distributed with the Mahout machine-learning library. We
played around with all our basic matchers as features, but ended up only using a quite small subset in the final
version. Only 6 features were used:
\begin{itemize}
    \item Lemma Matching
    \item IDF Lemma Matching
    \item Lemma+POS Matching
    \item BleuScore
    \item Polarity
    \item WordNet Synonym Matching
\end{itemize}

For all features for which we had matchers (all except Polarity) we used the matchers estimate whether a text/hypothesis
pair was entailing or not as value, as this is already a convinient value between 0 and 1. All matchers except for
BleuScore already contain the Polarity measurement already as a sort of "deatch" criteria, where we set the estimate to
0 if the polarity doesn't match. We still use it as seperate feature, by setting the value to 1 if it matches and 0 if
it doesn't, this proved helpful (without polarity feature, the result was ~62.6\%).




\section{Results}

- TreeDistMatcher combined with cost function

\begin{center}
\begin{tabular}{ | c | c | c | c | }
	\hline
	technique 				& parameters				& threshold 	& correct	\\ 
	\hline
	MahoutMatcher			& 10 fold cross-validation	& avg. 0.5225	& 65.625\%	\\
							&							& [0.475-0.55]	& [57.5\%-75\%] \\
	\hline
	SynonymMatching			&							& 0.675			& 63.375\%	\\
	LemmaMatching			&							& 0.675			& 63.25\%	\\
	LemmaAndPosMatching 	&							& 0.625			& 62.875\%	\\	
  	LexicalMatching			& 							& 0.575			& 62\%		\\
	BleuScoreMatching		& depth=2					& 0.425			& 61.875\%	\\
	IDFLexicalMatching		&							& 0.325			& 61.625\%	\\
	IDFLemmaMatching		&							& 0.5			& 61.375\%	\\
	BleuScoreMatching		& depth=4					& 0.5			& 61.25\%	\\	
	BleuScoreMatching		& depth=3					& 0.425			& 61.125\%	\\	
  	BleuScoreMatching		& arithm. mean, depth=2		& 0.075			& 60.875\%	\\	
	LinSimilarityMatching	&							& 0.35			& 60.5\%	\\
	WordNetDistanceMatching	& 							& 0.775			& 60.25\%	\\
	TreeDistMatcher			& costs: WeightedLemma		& 0.725			& 57.125\%	\\
	TreeDistMatcher			& costs: WeightedIDF		& 0.4			& 55.75\%	\\
	BleuScoreMatching		& arithm. mean, depth=3		& 0.05			& 55.375\%	\\	
	BleuScoreMatching		& arithm. mean, depth=4  	& 0.05			& 51.875\%	\\	
	TreeDistMatcher			& costs: FreeDeletion		& 0.05			& 51.7\%	\\
	\hline
\end{tabular}
\end{center}	


- very low threshold on some
- consider seperately because of cross validation
\section{Future plans}
Polarity measures could be implemented as a sort of filter around normal matchers, which can be added on demand.
Currently the polarity is used in all \textit{BasicMatchers}.
There are probably a lot of interesting possibilities with the WordNet API which have not been explored yet. Combining
these with a more tuned version of the Mahout learning algorithms might lead to better results in the future.

The code could be better documented using javadoc. We think that if the code is properly documented it could be released as a generic RTE java library that can be used by
others for their work. All code can be found at: \url{https://github.com/nihathrael/java-rte}



\end{document}
